\begin{ex}
Trong phép tính $32 + 15 = 47$, số 47 được gọi là gì? Tính tổng của 32 và 15.
\shortans{47}
\loigiai{
Tổng của hai số là 47.
}
\end{ex}

\begin{ex}
Trong phép tính $25 + 40 = 65$, số 25 được gọi là số hạng thứ nhất. Số hạng thứ hai bằng bao nhiêu?
\shortans{40}
\loigiai{
Số hạng thứ hai là 40.
}
\end{ex}

\begin{ex}
Trong phép tính $58 - 23 = 35$, số bị trừ bằng bao nhiêu?
\shortans{58}
\loigiai{
Số bị trừ là 58.
}
\end{ex}

\begin{ex}
Trong phép tính $76 - 34 = 42$, số trừ bằng bao nhiêu?
\shortans{34}
\loigiai{
Số trừ là 34.
}
\end{ex}

\begin{ex}
Trong phép tính $89 - 50 = 39$, hiệu bằng bao nhiêu?
\shortans{39}
\loigiai{
Hiệu là 39.
}
\end{ex}

\begin{ex}
Biết hai số hạng là 24 và 35. Tính tổng của hai số đó.
\shortans{59}
\loigiai{
Tổng: $24 + 35 = 59$.
}
\end{ex}

\begin{ex}
Biết số bị trừ là 67, số trừ là 25. Tính hiệu của phép trừ.
\shortans{42}
\loigiai{
Hiệu: $67 - 25 = 42$.
}
\end{ex}

\begin{ex}
Biết tổng bằng 58, một số hạng bằng 20. Tìm số hạng còn lại.
\shortans{38}
\loigiai{
Số hạng còn lại: $58 - 20 = 38$.
}
\end{ex}

\begin{ex}
Biết tổng bằng 76, một số hạng bằng 32. Tìm số hạng còn lại.
\shortans{44}
\loigiai{
Số hạng còn lại: $76 - 32 = 44$.
}
\end{ex}

\begin{ex}
Biết số bị trừ là 84, hiệu là 31. Tìm số trừ.
\shortans{53}
\loigiai{
Số trừ: $84 - 31 = 53$.
}
\end{ex}

\begin{ex}
Biết số trừ là 26, hiệu là 43. Tìm số bị trừ.
\shortans{69}
\loigiai{
Số bị trừ: $43 + 26 = 69$.
}
\end{ex}

\begin{ex}
Nếu số hạng thứ nhất tăng thêm 5 đơn vị và giữ nguyên số hạng thứ hai thì tổng tăng thêm bao nhiêu đơn vị?
\shortans{5}
\loigiai{
Tổng tăng thêm đúng 5 đơn vị.
}
\end{ex}

\begin{ex}
Nếu số bị trừ tăng thêm 8 đơn vị và giữ nguyên số trừ thì hiệu tăng thêm bao nhiêu đơn vị?
\shortans{8}
\loigiai{
Hiệu tăng thêm 8 đơn vị.
}
\end{ex}

\begin{ex}
Nếu giữ nguyên số bị trừ và số trừ tăng thêm 4 đơn vị thì hiệu giảm đi bao nhiêu đơn vị?
\shortans{4}
\loigiai{
Hiệu giảm đi 4 đơn vị.
}
\end{ex}

\begin{ex}
Tổng của hai số bằng 60, nếu một số hạng bằng 60 thì số hạng kia bằng bao nhiêu?
\shortans{0}
\loigiai{
Số hạng kia là: $60 - 60 = 0$.
}
\end{ex}

\begin{ex}
Hiệu của hai số bằng 0, biết số bị trừ là 45. Số trừ bằng bao nhiêu?
\shortans{45}
\loigiai{
Hai số bằng nhau thì hiệu bằng 0, nên số trừ là 45.
}
\end{ex}

\begin{ex}
Hiệu của hai số bằng 38, biết số trừ là 0. Số bị trừ bằng bao nhiêu?
\shortans{38}
\loigiai{
Số bị trừ: $38 + 0 = 38$.
}
\end{ex}

\begin{ex}
Tìm số bị trừ biết số trừ là số bé nhất có hai chữ số và hiệu là 52.
\shortans{62}
\loigiai{
Số trừ là 10. Số bị trừ: $52 + 10 = 62$.
}
\end{ex}

\begin{ex}
Tìm số hạng thứ hai biết tổng là số lớn nhất có hai chữ số và số hạng thứ nhất là 70.
\shortans{29}
\loigiai{
Tổng là 99. Số hạng thứ hai: $99 - 70 = 29$.
}
\end{ex}

\begin{ex}
Tìm hiệu biết số bị trừ là số tròn chục lớn nhất có hai chữ số và số trừ là 30.
\shortans{60}
\loigiai{
Số bị trừ là 90. Hiệu: $90 - 30 = 60$.
}
\end{ex}

\begin{ex}
Hai số có tổng bằng 88. Nếu số hạng thứ nhất là 55 thì số hạng thứ hai là bao nhiêu?
\shortans{33}
\loigiai{
Số hạng thứ hai: $88 - 55 = 33$.
}
\end{ex}

\begin{ex}
Tìm số trừ biết số bị trừ là 79 và hiệu là 46.
\shortans{33}
\loigiai{
Số trừ: $79 - 46 = 33$.
}
\end{ex}

\begin{ex}
Trong một phép cộng có hai số hạng bằng nhau và tổng bằng 18. Số hạng đó bằng bao nhiêu?
\shortans{9}
\loigiai{
Vì $9 + 9 = 18$ nên số hạng là 9.
}
\end{ex}

\begin{ex}
Trong một phép cộng có hai số hạng bằng nhau và tổng bằng 16. Số hạng đó bằng bao nhiêu?
\shortans{8}
\loigiai{
Vì $8 + 8 = 16$ nên số hạng là 8.
}
\end{ex}

\begin{ex}
Một phép trừ có số bị trừ bằng 50, số trừ bằng 50. Hiệu bằng bao nhiêu?
\shortans{0}
\loigiai{
Hiệu là: $50 - 50 = 0$.
}
\end{ex}
